% ViT on CIFAR-10 -- IEEE two-column report (4-6 pages)
% Build: pdflatex ViT_CIFAR10_IEEE_Report  (run twice).
% The numbers/tables come from ../results via `python analysis.py` (see the notebook, section 7):
%   generated/results_macros.tex, generated/tab_*.tex   and   ../figures/*.png
% Red text = a sentence that depends on YOUR measured numbers. Read it, edit it if the numbers
% disagree, then switch \drafttrue to \draftfalse below to turn the red off.
\documentclass[conference]{IEEEtran}
\IEEEoverridecommandlockouts
\usepackage{cite}
\usepackage{amsmath,amssymb}
\usepackage{graphicx}
\usepackage{booktabs}
\usepackage{xcolor}
\usepackage{url}
\usepackage[hidelinks]{hyperref}
\graphicspath{{../figures/}}

\newif\ifdraft \drafttrue
\newcommand{\verify}[1]{\ifdraft{\color{red}#1}\else#1\fi}

% ---- defaults, overwritten by generated/results_macros.tex ------------------------------
\newcommand{\accVitTen}{--}\newcommand{\accVitTwentyFive}{--}\newcommand{\accVitFifty}{--}\newcommand{\accVitFull}{--}
\newcommand{\accResTen}{--}\newcommand{\accResTwentyFive}{--}\newcommand{\accResFifty}{--}\newcommand{\accResFull}{--}
\newcommand{\gapTen}{--}\newcommand{\gapTwentyFive}{--}\newcommand{\gapFifty}{--}\newcommand{\gapFull}{--}
\newcommand{\paramsVit}{--}\newcommand{\paramsRes}{--}\newcommand{\epochsVit}{--}\newcommand{\epochsRes}{--}
\newcommand{\minutesVit}{--}\newcommand{\minutesRes}{--}
\newcommand{\accAblNone}{--}\newcommand{\accAblLearned}{--}\newcommand{\accAblSincos}{--}\newcommand{\accAblMean}{--}
\newcommand{\accFtVit}{--}\newcommand{\accFtRes}{--}\newcommand{\epochsFt}{--}
\newcommand{\topConfVit}{--}\newcommand{\topConfRes}{--}
\newcommand{\worstClassVit}{--}\newcommand{\worstAccVit}{--}\newcommand{\bestClassVit}{--}\newcommand{\bestAccVit}{--}
\newcommand{\worstClassRes}{--}\newcommand{\worstAccRes}{--}
\IfFileExists{generated/results_macros.tex}{\input{generated/results_macros.tex}}{}
\newcommand{\gentab}[1]{\IfFileExists{generated/#1.tex}{\input{generated/#1.tex}}{\emph{(table missing: run the notebook)}}}
\newcommand{\figbox}[2]{\IfFileExists{../figures/#1.png}{\includegraphics[width=#2]{#1}}{\fbox{\parbox{#2}{\centering\vspace{2.5em}missing figure: \detokenize{#1}\vspace{2.5em}}}}}

\title{Do Vision Transformers Need Convolutional Inductive Bias? A Reduced-Scale Reproduction of ``An Image is Worth 16$\times$16 Words'' on CIFAR-10}

\author{
\IEEEauthorblockN{Aditya Kumar Jha}
\IEEEauthorblockA{B.Tech CSE, RGIPT --- Course 24CS3006, Deep Learning 2026 (paper project DL-065)\\
Email: [add institutional email if required]}
}

\begin{document}
\maketitle

\begin{abstract}
The Vision Transformer (ViT) classifies an image by feeding a sequence of linearly embedded patches to a standard Transformer encoder. Its central claim is that this architecture, which lacks the locality and translation-equivariance priors of convolutions, matches or beats strong CNNs only when it is pre-trained on large datasets. I re-implement ViT from scratch in PyTorch (patch embedding, CLS token, position embeddings, hand-written multi-head self-attention, pre-LN encoder) and test this claim at small scale on CIFAR-10 on a single Colab T4 GPU. Experiments: (i) ViT vs.\ a CIFAR-adapted ResNet-18, both trained from scratch on 10--100\,\% of the training set; (ii) the same comparison after fine-tuning ImageNet-pre-trained ViT-B/16 and ResNet-50; (iii) an ablation of the position embedding (none, 1-D learned, 2-D sin-cos); and (iv) attention-rollout visualisations and a confusion analysis. \verify{With the full training set the ViT reaches \accVitFull\,\% test accuracy against \accResFull\,\% for ResNet-18; the ResNet's lead is \gapTen{} points at 10\,\% of the data and \gapFull{} at 100\,\%. Fine-tuned ViT-B/16 reaches \accFtVit\,\% versus \accFtRes\,\% for ResNet-50 (paper: 98.95\,\% for ImageNet-21k pre-training). Removing position embeddings gives \accAblNone\,\% versus \accAblLearned\,\% with learned embeddings.}
\end{abstract}

\begin{IEEEkeywords}
Vision Transformer, inductive bias, CIFAR-10, attention rollout, transfer learning, ablation.
\end{IEEEkeywords}

\section{Introduction and Motivation}
Convolutional networks build in two useful assumptions about images: \emph{locality} (nearby pixels interact) and \emph{translation equivariance} (the same filter is applied everywhere). Dosovitskiy \emph{et al.}\ \cite{dosovitskiy2021} showed that a pure Transformer \cite{vaswani2017}, applied to $16\times16$ image patches with almost no image-specific structure, reaches excellent accuracy after large-scale pre-training. They also report that, trained on a mid-sized dataset such as ImageNet without strong regularisation, ViT is a few points \emph{worse} than comparable ResNets; the lack of inductive bias must be paid for with data.

This project asks whether that trend is visible at a scale a student can reproduce. The questions are: \textbf{Q1} how does a from-scratch ViT compare with a ResNet as the amount of training data shrinks? \textbf{Q2} does pre-training on a large dataset close the gap, and how do my numbers compare with the paper's? \textbf{Q3} which ViT component matters most (here: the position embedding)? \textbf{Q4} what does the model attend to, and where does it fail? The paper's JFT-300M/ImageNet-21k pre-training is far outside the compute budget, so the reproduction is explicitly \emph{reduced-scale}: what changed is stated in Sec.~\ref{sec:setup}.

\section{Method}
\subsection{Vision Transformer}
An image $\mathbf{x}\in\mathbb{R}^{H\times W\times C}$ is cut into $N=HW/P^2$ non-overlapping patches $\mathbf{x}_p^i\in\mathbb{R}^{P^2C}$ which are flattened and projected by one shared matrix $\mathbf{E}\in\mathbb{R}^{(P^2C)\times D}$. A learnable CLS token and position embeddings $\mathbf{E}_{pos}\in\mathbb{R}^{(N+1)\times D}$ are added:
\begin{equation}
\mathbf{z}_0=[\mathbf{x}_{cls};\,\mathbf{x}_p^1\mathbf{E};\dots;\mathbf{x}_p^N\mathbf{E}]+\mathbf{E}_{pos}.
\end{equation}
Each of the $L$ pre-LN encoder blocks computes
\begin{align}
\mathbf{z}'_\ell&=\mathrm{MSA}(\mathrm{LN}(\mathbf{z}_{\ell-1}))+\mathbf{z}_{\ell-1},\\
\mathbf{z}_\ell&=\mathrm{MLP}(\mathrm{LN}(\mathbf{z}'_\ell))+\mathbf{z}'_\ell,
\end{align}
and the class logits are $\mathbf{y}=\mathbf{W}\,\mathrm{LN}(\mathbf{z}_L^{0})$ (the CLS row). Self-attention for head $h$ is
\begin{equation}
\mathrm{head}_h=\mathrm{softmax}\!\Big(\frac{\mathbf{Q}_h\mathbf{K}_h^{\top}}{\sqrt{d_h}}\Big)\mathbf{V}_h,\quad d_h=D/\text{heads}.
\label{eq:attn}
\end{equation}
The factor $1/\sqrt{d_h}$ keeps the variance of the dot products near 1 (for unit-variance, independent query and key entries, $\mathrm{Var}(\mathbf{q}\!\cdot\!\mathbf{k})=d_h$), which stops the softmax from saturating and its gradients from vanishing. The patch projection is implemented as a \texttt{Conv2d} with kernel and stride $P$, which is exactly equivalent to flatten-then-\texttt{Linear} (verified by a unit test in the repository); it has $(P^2C)D+D$ parameters.

\subsection{Position-embedding variants}
Besides the paper's learned 1-D embedding, I implement a fixed 2-D sine-cosine embedding (half of the channels encode the patch row, half the column, no trainable parameters) and a ``none'' variant in which the encoder sees an unordered bag of patches.

\subsection{Attention rollout}
Raw attention of the last layer ignores how information is mixed in earlier layers. Attention rollout \cite{abnar2020} propagates it through the network: with $\bar{\mathbf{A}}_\ell$ the head-averaged attention of layer $\ell$ and the identity added for the residual connection,
\begin{equation}
\hat{\mathbf{A}}_\ell=\mathrm{rownorm}\big(0.5\,\bar{\mathbf{A}}_\ell+0.5\,\mathbf{I}\big),\qquad
\mathbf{R}=\hat{\mathbf{A}}_L\cdots\hat{\mathbf{A}}_1 .
\end{equation}
The CLS row of $\mathbf{R}$, without the CLS column, is reshaped to the $8\times8$ patch grid and bilinearly up-sampled onto the image. Rollout is a diagnostic of information flow, not a causal explanation.

\section{Experimental Setup}\label{sec:setup}
\textbf{Data.} CIFAR-10 \cite{krizhevsky2009}: a fixed 45\,000/5\,000 train/validation split of the training images (seed 42); the 10\,000 test images are touched once, on the checkpoint with the best validation accuracy. Subsets for the data-fraction study are nested (10\,\%\,$\subset$\,25\,\%\,$\subset\cdots$).

\textbf{Models.} ViT: $32\times32$ input, patch size 4 ($N{=}64$ tokens $+$ CLS), $D{=}192$, depth 8, 6 heads, MLP width 768, \paramsVit{} parameters (a ViT-Tiny-sized model; the paper's ViT-B/16 has 86\,M). ResNet-18 \cite{he2016} with the usual CIFAR stem (3$\times$3 stride-1 convolution, no max-pool), \paramsRes{} parameters, trained from scratch.

\textbf{Training recipe (identical for both from-scratch models).} Random crop (reflect-padding 4), horizontal flip, RandAugment \cite{cubuk2020}, mixup \cite{zhang2018} / CutMix \cite{yun2019} (50/50), label smoothing 0.1 \cite{szegedy2016}, linear warm-up and cosine decay, batch size 128, fp16 autocast on a Colab T4. ViT: AdamW \cite{loshchilov2019}, lr $10^{-3}$, weight decay 0.05, gradient clipping 1; ResNet: Nesterov SGD, lr 0.1, weight decay $5{\times}10^{-4}$. \verify{The 100\,\% runs use \epochsVit{} epochs; smaller subsets are trained for proportionally more epochs so that every model receives the same number of updates (epochs are listed in Table~\ref{tab:sweep}), separating the effect of data quantity from that of training length.} Everything is seeded; the repository README lists the exact commands.

\textbf{Fine-tuning.} Pre-trained \texttt{timm} \cite{wightman2019} ViT-B/16 and ResNet-50 get a new 10-way head, inputs are up-sampled from $32$ to $224$ on the GPU, and training uses SGD with momentum (lr 0.01, cosine, gradient clipping) for \epochsFt{} epochs. \verify{The ViT-B/16 checkpoint is \texttt{timm}'s default \texttt{vit\_base\_patch16\_224} (AugReg training \cite{steiner2022}, ImageNet-21k pre-training, ImageNet-1k fine-tuning), not the original paper checkpoint; check the tag printed by your timm version.}

\textbf{What differs from the paper.} (a) CIFAR-10 instead of ImageNet/JFT pre-training for the from-scratch study; (b) a ViT-Tiny-sized model; (c) tens of epochs instead of hundreds; (d) \texttt{timm} weights and a few fine-tuning epochs instead of the paper's fine-tuning schedule; (e) one seed per configuration unless a $\pm$ is shown.

\begin{table}[t]
\centering\small
\caption{Training hyper-parameters.}
\label{tab:hp}
\begin{tabular}{@{}lccc@{}}
\toprule
 & ViT-Tiny & ResNet-18 & Fine-tune \\
\midrule
Optimiser & AdamW & Nesterov SGD & SGD (mom.\ 0.9) \\
Peak lr & $10^{-3}$ & 0.1 & 0.01 \\
Weight decay & 0.05 & $5{\times}10^{-4}$ & 0 \\
Batch size & 128 & 128 & 64 \\
Schedule & \multicolumn{2}{c}{warm-up + cosine} & warm-up + cosine \\
Label smoothing & 0.1 & 0.1 & 0 \\
Mixup/CutMix prob. & \verify{0.5--0.8} & \verify{0.5--0.8} & 0 \\
Input & $32^2$ & $32^2$ & $224^2$ (up-sampled) \\
\bottomrule
\end{tabular}
\end{table}

\section{Results}
\subsection{Q1: from-scratch ViT vs.\ ResNet-18 across data sizes}
\begin{table}[t]
\centering\small
\caption{Test accuracy (\%) on CIFAR-10, from scratch, versus training-set size. Gap = ResNet $-$ ViT in percentage points.}
\label{tab:sweep}
\gentab{tab_sweep}
\end{table}
\begin{figure}[t]
\centering\figbox{fraction_sweep}{0.85\columnwidth}
\caption{Data-scale study in miniature (compare the paper's Fig.~3--4): accuracy versus fraction of the training set.}
\label{fig:sweep}
\end{figure}
\begin{figure}[t]
\centering\figbox{curves_full_data}{\columnwidth}
\caption{Validation loss and accuracy of the two from-scratch models trained on the full training set.}
\label{fig:curves}
\end{figure}
Table~\ref{tab:sweep} and Fig.~\ref{fig:sweep} give the main result (training curves: Fig.~\ref{fig:curves}). \verify{ResNet-18 is ahead of the ViT at every data size, by \gapFull{} points with all 45\,000 images and \gapTen{} points with 10\,\% (4\,500 images). The ViT loses more accuracy than the ResNet when data are removed, which is the qualitative trend the paper describes: without convolutional priors the transformer must learn locality from data.} Both absolute numbers are limited by the short schedule and by the small image size, which gives the ViT only 64 patch tokens.

\subsection{Q2: pre-training and comparison with the paper}
\begin{table}[t]
\centering\small
\caption{Fine-tuned on CIFAR-10 (this work) versus the paper's Table~5 for ViT-B/16.}
\label{tab:ft}
\gentab{tab_finetune}
\end{table}
\verify{Pre-training removes the gap: ViT-B/16 reaches \accFtVit\,\% and ResNet-50 \accFtRes\,\% after only \epochsFt{} epochs (Table~\ref{tab:ft}), compared with \accVitFull\,\% and \accResFull\,\% when ViT-Tiny and ResNet-18 are trained from scratch. My ViT-B/16 figure should be read against the paper's 98.95\,\% for ImageNet-21k pre-training and 98.13\,\% for ImageNet-1k pre-training; the remaining difference is explained by the different checkpoint, resolution handling and the much shorter fine-tuning schedule listed in Sec.~\ref{sec:setup}.} The paper reports no from-scratch CIFAR-10 number, so the from-scratch rows are my own addition.

\subsection{Q3: position-embedding ablation}
\begin{table}[t]
\centering\small
\caption{Position-embedding ablation (ViT, full training set, same recipe).}
\label{tab:abl}
\gentab{tab_ablation}
\end{table}
Following the paper's own ablation (its Appendix D.4), which found that \emph{having} position information matters while \emph{how} it is encoded matters little, Table~\ref{tab:abl} compares no position embedding, 1-D learned and 2-D sin-cos. \verify{Without position embeddings the accuracy is \accAblNone\,\%, against \accAblLearned\,\% (learned) and \accAblSincos\,\% (sin-cos). Differences between the two position encodings are small relative to the gap to ``none''; with one seed per variant, differences under about one point should not be over-interpreted.} Without positions the model is permutation-invariant over patches and cannot use spatial layout beyond what is present inside a single $4\times4$ patch.

\section{Qualitative Analysis and Failure Cases}
\begin{figure}[t]
\centering\figbox{rollout_grid}{0.78\columnwidth}
\caption{Attention rollout of the full-data ViT. Rows: correctly classified images, their rollout overlays, misclassified images (T = true, P = predicted) and overlays.}
\label{fig:rollout}
\end{figure}
Fig.~\ref{fig:rollout} shows rollout maps for four correctly and four incorrectly classified test images. \verify{[Describe what you actually see: e.g.\ whether the CLS token concentrates on the object or on the background, and whether misclassified images show diffuse attention or attention on a distracting region. Do not claim object-centred attention unless the figure shows it.]} Because rollout multiplies averaged attention matrices it is a coarse diagnostic and says nothing about causal importance.

\begin{figure}[t]
\centering\figbox{confusion_vit}{0.8\columnwidth}
\caption{Row-normalised confusion matrix of the full-data ViT on the test set.}
\label{fig:cm}
\end{figure}
Fig.~\ref{fig:cm} shows the ViT's errors. \verify{The most frequent confusions are \topConfVit{} (true$\to$predicted, number of test images), while the easiest and hardest classes are \bestClassVit{} (\bestAccVit\,\%) and \worstClassVit{} (\worstAccVit\,\%). For ResNet-18 the most frequent confusions are \topConfRes{}. [Add one sentence on why, e.g.\ whether the confused classes are visually similar at $32\times32$ resolution, only if the matrix supports it.]}

\textbf{What did not work / caveats.} A first version of this project (15 epochs, crop and flip only, dropout 0.1, $D{=}128$, 1.2\,M parameters) reached only 72.0\,\% for the ViT and 89.4\,\% for ResNet-18; its dropout ablation (75.0\,\% without dropout) used different seeds per run and is not repeated here. Without strong augmentation a from-scratch ViT overfits CIFAR-10 quickly, which is why the final recipe adds RandAugment, mixup/CutMix, label smoothing and warm-up.

\section{Discussion and Limitations}
\verify{The experiments support the paper's qualitative message in miniature: the transformer is competitive only with a lot of data or with pre-training, while a convolutional network with its built-in priors is data-efficient.} The limitations are real: \verify{one seed per configuration (differences of a few tenths of a point are within noise)}, short schedules, a small ViT, $32\times32$ inputs with only 64 tokens, a hyper-parameter recipe that was tuned lightly and not per data fraction, and checkpoint selection on a 5\,000-image validation set. The AugReg-trained \texttt{timm} checkpoint is not the paper's. These choices were made for a free-GPU budget, and the absolute numbers should not be compared with state-of-the-art CIFAR-10 results.

\section{Conclusion}
A ViT was written from scratch (patch embedding, CLS token, learned/sin-cos/no position embedding, manual multi-head attention, pre-LN blocks, rollout) and compared with ResNets on CIFAR-10 in three regimes: from scratch with shrinking data, fine-tuned from pre-trained weights, and with position embeddings ablated. \verify{ViT trails ResNet-18 from scratch and more so with less data, pre-training closes the gap, and position information is the component whose removal costs the most.} The code, seeds, results and a Colab notebook are in the project repository \url{https://github.com/aditya-kumar04/ViT-CIFAR10}. The ViT paper's official code was not used; only the paper (and \cite{abnar2020} for rollout) was consulted, and \texttt{timm} supplies the pre-trained backbones.

\begin{thebibliography}{00}
\bibitem{dosovitskiy2021}
A. Dosovitskiy \emph{et al.}, ``An image is worth 16x16 words: Transformers for image recognition at scale,'' in \emph{Proc. ICLR}, 2021.
\bibitem{vaswani2017}
A. Vaswani \emph{et al.}, ``Attention is all you need,'' in \emph{Proc. NeurIPS}, 2017.
\bibitem{he2016}
K. He, X. Zhang, S. Ren, and J. Sun, ``Deep residual learning for image recognition,'' in \emph{Proc. CVPR}, 2016, pp. 770--778.
\bibitem{krizhevsky2009}
A. Krizhevsky, ``Learning multiple layers of features from tiny images,'' Univ.\ Toronto, Tech.\ Rep., 2009.
\bibitem{abnar2020}
S. Abnar and W. Zuidema, ``Quantifying attention flow in Transformers,'' in \emph{Proc. ACL}, 2020, pp. 4190--4197.
\bibitem{cubuk2020}
E. D. Cubuk, B. Zoph, J. Shlens, and Q. V. Le, ``RandAugment: Practical automated data augmentation with a reduced search space,'' in \emph{Proc. CVPR Workshops}, 2020.
\bibitem{zhang2018}
H. Zhang, M. Cisse, Y. N. Dauphin, and D. Lopez-Paz, ``mixup: Beyond empirical risk minimization,'' in \emph{Proc. ICLR}, 2018.
\bibitem{yun2019}
S. Yun \emph{et al.}, ``CutMix: Regularization strategy to train strong classifiers with localizable features,'' in \emph{Proc. ICCV}, 2019.
\bibitem{szegedy2016}
C. Szegedy, V. Vanhoucke, S. Ioffe, J. Shlens, and Z. Wojna, ``Rethinking the Inception architecture for computer vision,'' in \emph{Proc. CVPR}, 2016, pp. 2818--2826.
\bibitem{loshchilov2019}
I. Loshchilov and F. Hutter, ``Decoupled weight decay regularization,'' in \emph{Proc. ICLR}, 2019.
\bibitem{steiner2022}
A. Steiner \emph{et al.}, ``How to train your ViT? Data, augmentation, and regularization in Vision Transformers,'' \emph{Trans.\ Mach.\ Learn.\ Res.}, 2022.
\bibitem{wightman2019}
R. Wightman, ``PyTorch image models (timm),'' GitHub repository, 2019. [Online]. Available: \url{https://github.com/huggingface/pytorch-image-models}
\end{thebibliography}
\end{document}
